\documentclass[11pt,a4paper]{article}
\usepackage{cmap}
\usepackage[utf8]{inputenc}
\usepackage[T1]{fontenc}
\usepackage[margin=0.78in]{geometry}
\usepackage{amsmath,amssymb}
\usepackage{booktabs}
\usepackage{enumitem}
\usepackage{xcolor}
\usepackage{colortbl}
\usepackage{hyperref}
\usepackage{setspace}
\usepackage{fancyhdr}
\usepackage{lastpage}
\hypersetup{colorlinks=true,linkcolor=blue!60!black,
  pdftitle={Reviewer Roadmap -- Spectral Geometry on T5 v46},
  pdfauthor={Martin Wolfgang Le Borgne},
  pdfsubject={m_p/m_e = 6 pi^5 (1 + alpha^2/sqrt(8)) = 1836.15268; 20 parameter-free empirical consequences; Working Hypothesis M: unique self-consistent fixed point},
  pdfkeywords={SFST, reviewer roadmap, falsifiable predictions, spectral geometry}}
\setlength{\parskip}{0.3em}
\setlength{\parindent}{0pt}
\pagestyle{fancy}\fancyhf{}
\lhead{SFST v46 --- Reviewer Roadmap}\rhead{\thepage/\pageref{LastPage}}
\setlength{\headheight}{14pt}

\definecolor{tier1}{HTML}{D4EDDA}
\definecolor{tier2}{HTML}{FFF3CD}
\definecolor{tier3}{HTML}{F8D7DA}

\title{\vspace{-2em}\textbf{Reviewer Roadmap}\\[0.3em]
\large Spectral Geometry on $T^5$ with SU(3): Algebraic Identities and Their Physical Consequences\\[0.2em]
\normalsize v46 --- March 2026}
\author{Martin Wolfgang Le Borgne}
\date{}

\begin{document}
\maketitle\thispagestyle{fancy}\vspace{-2em}

\section*{What the paper claims}

One axiom (the matching map, Axiom~M) plus the spectral geometry of
$T^5$ with SU(3) gauge field yields \textbf{20 parameter-free empirical consequences},
including three high-precision numerical values:
\begin{center}
\renewcommand{\arraystretch}{1.1}
\begin{tabular}{lll}
\toprule
\textbf{Prediction} & \textbf{Formula} & \textbf{Accuracy}\\
\midrule
V1: proton--electron mass ratio & $m_p/m_e = 6\pi^5(1+\alpha^2/\!\sqrt{8})$ & 0.002\,ppm\\
V3: fine-structure constant & $-2\ln\alpha = \pi^2-4\alpha+c_2\alpha^2$ & 0.009\,ppm\\
V2: neutron--proton splitting & $\Delta m/m_e = 8/\pi-2\alpha$ & 354\,ppm\\
\bottomrule
\end{tabular}
\end{center}

\section*{Proof architecture: what is proven and what is not}

\begin{center}
\renewcommand{\arraystretch}{1.15}\small
\begin{tabular}{>{\raggedright}p{5.5cm}cp{5.5cm}}
\toprule
\textbf{Result} & \textbf{Tier} & \textbf{Location / Script}\\
\midrule
\rowcolor{tier1}
$d=5$ uniqueness (4 conditions, unique intersection)
  & 1 & Supp.\ Z1; \texttt{d5\_uniqueness.py}\\
\rowcolor{tier1}
Ramond BC uniqueness (permutation + sign + Hosotani)
  & 1 & Supp.\ Z2; \texttt{ramond\_bc\_proof.py}\\
\rowcolor{tier1}
Hosotani global stability ($H=32.3\times I_5$)
  & 1 & Supp.\ Z3; \texttt{hosotani\_global\_stability.py}\\
\rowcolor{tier1}
$Q_u\!=\!2/3$, $Q_d\!=\!-1/3$ from $Q_p\!=\!1$, $Q_n\!=\!0$
  & 1 & Supp.\ Z4; \texttt{quark\_charge\_derivation.py}\\
\rowcolor{tier1}
$c_2=\tfrac52\ln 2-\tfrac38$ unique (sensitivity test)
  & 1 & Supp.\ Z5; \texttt{c2\_schema\_independence.py}\\
\rowcolor{tier1}
$1/\!\sqrt{8}$ from 3 routes (scheme-independent $\alpha_s$)
  & 1 & Supp.\ Z, Z6; \texttt{smooth\_cutoff\_closing\_gap.py}\\
\rowcolor{tier1}
$R$-independence ($6\pi^5$ for all $R$; no T-duality needed)
  & 1 & Supp.\ Z7; \texttt{R\_half\_independence.py}\\
\rowcolor{tier1}
$\alpha^1$-cancellation ($\sum Q^2=1$)
  & 1 & Supp.\ Z; \texttt{trace\_computation.py}\\
\rowcolor{tier1}
Weyl identity ($6\pi^5$ exact)
  & 1 & Paper Thm.\ (Weyl identity)\\
\rowcolor{tier1}
Epstein zeta uniqueness (50-digit)
  & 1 & \texttt{epstein\_zeta\_direct.py}\\
\midrule
\rowcolor{tier2}
$c_3\approx 0.051$ (numerical, no closed form)
  & 2 & \texttt{c3\_derivation.py}\\
\rowcolor{tier2}
$\sin^2\theta_W(\Lambda)=3/8$ (qualitative)
  & 2--3 & Supp.\ Z8\\
\midrule
\rowcolor{tier1}
Toy model: matching is an \emph{identity} in $d=1,\ldots,5$
  & 1 & Supp.\ Z13; \texttt{toy\_model\_axiom\_m.py}\\
\rowcolor{tier1}
Triple verification of $\bar K=\pi^{5/2}$ (30+ digits, 3 methods)
  & 1 & Supp.\ Z10; \texttt{triple\_normalization\_check.py}\\
\rowcolor{tier2}
EW correction: $\delta_{\mathrm{EW}}\approx 3\,$ppb
  (W cancels; Z residual only)
  & 2 & Supp.\ Z11; \texttt{ew\_2loop\_upper\_bound.py}\\
\rowcolor{tier2}
Bootstrap closure: WH~M is the unique self-consistent fixed point
  & 1--2 & Supp.\ Z12; \texttt{bootstrap\_closure.py}\\
\midrule
\rowcolor{tier3}
\textbf{Working Hypothesis~M} (matching map: spectral data $\to$ mass ratio)
  & 3 & Paper \S 5; \textbf{structure proven, value tested by 20 predictions}\\
\rowcolor{tier3}
$\pi^7$-bridge ($6/\pi^2\to 6\pi^5$; equiv.\ to QCD mass gap on $T^5$)
  & 3 & Supp.\ Z13\\
\bottomrule
\end{tabular}\\[0.5em]
\colorbox{tier1}{\strut\small\ Tier~1: proven\ }\quad
\colorbox{tier2}{\strut\small\ Tier~2: closed with small residual\ }\quad
\colorbox{tier3}{\strut\small\ Tier~3: open\ }
\end{center}

\section*{Working Hypothesis~M --- why it is maximally constrained}

\textbf{Working Hypothesis~M} connects 5D spectral geometry (UV) to 4D hadron masses (IR).
It is supported by four independent lines of evidence:

\paragraph{1.\ Structural uniqueness (Z9).}
Five general principles (P1--P5) force the \emph{unique} form
$m_p/m_e = |W|\cdot[\bar K(\sigma^*)]^{n_p-n_e}$:

\begin{center}
\renewcommand{\arraystretch}{1.15}\small
\begin{tabular}{clp{6.8cm}}
\toprule
& \textbf{Principle} & \textbf{What it forces}\\
\midrule
P1 & Spectral origin & $\ln(m_p/m_e)$ linear in spectral data\\
P2 & Multiplicativity & data $=$ heat kernel $\Theta_3(1)^5$\\
P3 & $R$-independence & argument is $\sigma^*/R^2=1$ (self-dual point)\\
P4 & Topological quantization & exponent $n_p-n_e\in\mathbb{Z}$\\
P5 & Gauge covariance & overall factor $|W(\mathrm{SU}(3))|=6$\\
\bottomrule
\end{tabular}
\end{center}

\paragraph{2.\ Toy-model realization (Z13).}
On product tori $T^d$ ($d=1,\ldots,5$), the matching map is not a
hypothesis but an \emph{identity}: the spectral determinant ratio
gives $m_p/m_e = [\sin(\pi\theta)/\pi]^{d/2} = \pi^{-d/2} = 1/\bar K$
\emph{exactly}. The SFST adds two standard physics ingredients:
$|W|=6$ (color-singlet projection) and $\Delta n=2$ (instanton topology).

\paragraph{3.\ Bootstrap closure (Z12).}
Working Hypothesis~M is the \emph{unique} self-consistent fixed point:

\begin{center}
\renewcommand{\arraystretch}{1.1}\small
\begin{tabular}{lll}
\toprule
\textbf{Parameter} & \textbf{Allowed} & \textbf{Alternative}\\
\midrule
$|W|$ & $6$ (unique integer) & $5\to -17\%$; $7\to +17\%$\\
$\Delta n$ & $2$ (unique integer) & $1\to -94\%$; $3\to +1649\%$\\
$c_2$ & $[0.86,\,1.86]$ & Elizalde $1.358$ is central\\
$\delta_{\mathrm{EW}}$ & $[-8.3,\,+3.7]\,$ppb & estimate ${\approx}\,3\,$ppb $\checkmark$\\
\bottomrule
\end{tabular}
\end{center}

\paragraph{4.\ Triple normalization check (Z10).}
$\bar K = \pi^{5/2}$ verified to 30+ digits by three independent methods:
(i)~algebraic (Seeley--DeWitt),
(ii)~Chowla--Selberg lattice sum,
(iii)~theta function $+$ Poisson resummation.
Seeley--DeWitt coefficients $a_0,a_1$ cancel exactly in the ratio;
$a_2$ gives the finite $\alpha^2/\!\sqrt{8}$ correction.

\paragraph{The Lovelock analogy (Z9).}
In GR, Lovelock's theorem determines the \emph{structure} of Einstein's
equation from general principles; $\Lambda$ remains free.
WH~M is analogous: P1--P5 determine its form, the toy model confirms
its mechanism, and the bootstrap shows it is the unique fixed point.

\paragraph{Electroweak correction (Z11).}
The $W$-boson contribution \emph{cancels} in the ratio
(same $\mathrm{SU}(2)$ Casimir $C_2=3/4$ for quarks and electrons).
The surviving $Z$-boson correction is $\delta_{\mathrm{EW}} \approx 3\,$ppb
(range $0.5$--$50\,$ppb), consistent with the $+2.3\,$ppb residual
and below current experimental precision (${\sim}\,6\,$ppb).

\section*{Falsifiable tests}

\begin{center}
\renewcommand{\arraystretch}{1.15}\small
\begin{tabular}{lp{4cm}p{4.5cm}}
\toprule
\textbf{Prediction} & \textbf{Experiment} & \textbf{What kills SFST}\\
\midrule
V1: $m_p/m_e$ to $10^{-11}$
  & Penning trap (BASE, ALPHATRAP)
  & Deviation $>0.01\,$ppm at known $\alpha$\\
V3: $\alpha$ from $\pi$, $\ln 2$
  & Cs-133 interferometry
  & 2-loop value disagrees at $10^{-10}$\\
V18: $\theta_{\mathrm{QCD}}=0$
  & nEDM (PSI, SNS) at $10^{-28}\,e\!\cdot\!\mathrm{cm}$
  & Nonzero nEDM or axion detection\\
V20: Proton stable
  & Hyper-K, DUNE to $10^{35}\,$yr
  & \emph{Any} proton decay event\\
V21: $N_{\mathrm{gen}}=3$
  & LHC Run~3+
  & 4th-generation fermion discovery\\
\bottomrule
\end{tabular}
\end{center}

\section*{Key implications if confirmed}

\begin{enumerate}[leftmargin=1.5em,itemsep=0.2em]
\item $\alpha$ is computable $\Rightarrow$ 18 not 19 SM parameters; multiverse
  motivation for $\alpha$ collapses.
\item $\theta_{\mathrm{QCD}}=0$ without axion $\Rightarrow$ dark-matter
  searches must redirect from axion detectors (ADMX, MADMAX, CASPEr).
\item $d=5$ not $10$ or $11$ $\Rightarrow$ string theory in standard form
  is incomplete; the landscape problem would be significantly constrained.
\item Proton absolutely stable $\Rightarrow$ GUT-mediated proton decay channels are excluded;
  baryogenesis via leptogenesis.
\item $N_c=3$ is geometric $\Rightarrow$ the strong force is determined by
  the dimensionality of space, not by accident.
\end{enumerate}

\section*{Computational reproducibility}

All numerical claims are verified by Python scripts in the repository
(\texttt{scripts/} directory, 36 files). Each script runs in under 5\,min
on a standard laptop or Google Colab. Key dependencies: \texttt{mpmath}
(50-digit arithmetic), \texttt{numpy}, \texttt{scipy}.

\paragraph{Key scripts.}
\texttt{toy\_model\_axiom\_m.py} (identity in $d=1$--$5$);
\texttt{triple\_normalization\_check.py} (3 methods, 30 digits);
\texttt{bootstrap\_closure.py} (unique fixed point);
\texttt{ew\_2loop\_upper\_bound.py} ($W$ cancellation $+$ $Z$ bound);
\texttt{explicit\_1loop\_matching.py} (full derivation of $6\pi^5$).

\vfill

\end{document}
